\documentclass[troisieme]{fiche}
\metadonnees{14}{Un nouveau maître}{Bloc C — Décrire et comparer}

\begin{document}
\entete

\objectifs{%
\textbf{Langue} : les relatifs \emph{who}, \emph{which}, \emph{that}, \emph{whose},
\emph{where}.\par
\textbf{Texte} : Anna Sewell, \emph{Black Beauty} (1877), ch. \textsc{xxxiii}.\par
\textbf{Durée} : 1 h — exercices 20 min, texte et questions 25 min, version 15 min.}

%% ===================================================================
\exercice[12 min]{Choisir le relatif}

\rappel{\textbf{Rappel.} \emph{Who} pour les personnes, \emph{which} pour les choses,
\emph{that} pour les deux. \emph{Whose} = \og dont \fg{} (possession), \emph{where} = \og
où \fg.}

\consigne{Complétez.}

\begin{anglais}
\begin{enumerate}[label=\arabic*.,itemsep=2.5mm,leftmargin=*]
  \item Jerry is the man \trou{who} bought me at the fair.
  \item The cab \trou{which} he drove was his own.
  \item Polly is the woman \trou{whose} eyes are dark.
  \item This is the street \trou{where} the cab stand is.
  \item The horse \trou{that} they called Captain went to the war.
  \item The bread \trou{that} Dolly brought me was fresh.
  \item The men \trou{who} were sitting on their boxes were reading the newspaper.
  \item There was plenty of room in the stall, \trou{which} was clean and dry.
  \item I did not understand the word \trou{that} Harry meant.
  \item The lady \trou{who} spoke to us was rude.
\end{enumerate}
\end{anglais}

\note[\emph{That} ou \emph{which}]{\emph{That} remplace \emph{who} et \emph{which} quand la
relative est indispensable au sens (5, 6, 9). Mais après une virgule, la relative ajoute un
commentaire : \emph{that} y est impossible, il faut \emph{which} ou \emph{who} (8).}

\note[Quand le relatif peut disparaître]{S'il est complément d'objet, on peut l'omettre :
\emph{the bread Dolly brought me}, \emph{the word Harry meant} (6, 9). S'il est sujet, jamais :
\emph{the man who bought me} (1).}

%% ===================================================================
\exercice[8 min]{Relier deux phrases}

\rappel{\textbf{Rappel.} Le relatif se place juste après le nom qu'il reprend, et ce nom n'est
pas répété par un pronom.}

\consigne{Reliez les deux phrases par un pronom relatif.}

\begin{anglais}
\begin{enumerate}[label=\arabic*.,itemsep=3mm,leftmargin=*]
  \item This is the horse. It won the race.
    \lignes{1}
    \corr{This is the horse that won the race.}
  \item Jerry has a wife. Her name is Polly.
    \lignes{1}
    \corr{Jerry has a wife whose name is Polly.}
  \item I remember the day. We arrived in London on that day.
    \lignes{1}
    \corr{I remember the day when we arrived in London.}
  \item He showed me the stable. He keeps his horses there.
    \lignes{1}
    \corr{He showed me the stable where he keeps his horses.}
  \item The man was kind. I met him at the fair.
    \lignes{1}
    \corr{The man I met at the fair was kind. \emph{ou} \ldots{} that I met \ldots}
  \item Captain is an old horse. He went to the Crimean War.
    \lignes{1}
    \corr{Captain is an old horse that went to the Crimean War.}
\end{enumerate}
\end{anglais}

\note[La faute à ne pas faire]{Le pronom de la deuxième phrase disparaît avec la soudure :
\emph{the man I met at the fair} et non \emph{the man I met him at the fair}. Le français
commet rarement cette faute, l'anglais scolaire souvent.}

%% ===================================================================
\newpage
\section{Texte}

\noindent\small\textit{Le cheval qui raconte l'histoire a connu de mauvais maîtres. Il vient
d'être acheté dans une foire par un cocher de Londres.}\normalsize

\begin{texte}
Jeremiah Barker was my new master's name, but as every one called him Jerry, I shall do the
same. Polly, his wife, was just as good a \gl[a match]{match}{un parti, un assortiment} as a
man could have. She was a \gl[plump]{plump}{potelé}, \gl[trim]{trim}{soigné},
\gl[tidy]{tidy}{net, rangé} little woman, with smooth, dark hair, dark eyes, and a merry
little mouth. The boy was twelve years old, a tall, \gl[frank]{frank}{franc},
good-tempered \gl[a lad]{lad}{un garçon}; and little Dorothy (Dolly they called her) was her
mother over again, at eight years old. They were all wonderfully
\gl[fond of]{fond}{aimer beaucoup} of each other; I never knew such a happy, merry family
before or since. Jerry had a cab of his own, and two horses, which he drove and
\gl[to attend to]{attended to}{s'occuper de} himself. His other horse was a tall, white,
rather large-boned animal called ``Captain''. He was old now, but when he was young he must
have been splendid; he had still a proud way of holding his head and
\gl[to arch]{arching}{cambrer} his neck; in fact, he was a high-bred, fine-mannered, noble old
horse, every inch of him. He told me that in his early youth he went to the Crimean War; he
belonged to an officer in the \gl[the cavalry]{cavalry}{la cavalerie}, and used to lead the
\gl[a regiment]{regiment}{un régiment}. I will tell more of that hereafter.

The next morning, when I was well-\gl[to groom]{groomed}{panser}, Polly and Dolly came into
the yard to see me and make friends. Harry had been helping his father since the early
morning, and had stated his opinion that I should turn out a \gl[a regular brick]{regular
brick}{un brave type, un chic type}. Polly brought me a \gl[a slice]{slice}{une tranche} of
apple, and Dolly a piece of bread, and made as much of me as if I had been the ``Black
Beauty'' of \gl[olden]{olden time}{autrefois}. It was a great \gl[a treat]{treat}{un plaisir,
une gâterie} to be \gl[to pet]{petted}{caresser, choyer} again and talked to in a gentle
voice, and I let them see as well as I could that I wished to be friendly. Polly thought I was
very handsome, and a great deal too good for a cab, if it was not for the broken knees.

``Of course there's no one to tell us whose \gl[a fault]{fault}{une faute, la responsabilité}
that was,'' said Jerry, ``and as long as I don't know I shall give him the benefit of the
\gl[a doubt]{doubt}{un doute}. We'll call him `Jack', after the old one --- shall we, Polly?''

``Do,'' she said, ``for I like to keep a good name going.''
\end{texte}
\source{Anna Sewell, \emph{Black Beauty}, Londres, Jarrold, 1877, ch.~\textsc{xxxiii}.}

\partiel[15 min]{Questions}
\consigne{Répondez aux deux premières questions en français, aux deux dernières en anglais.}

\begin{enumerate}[label=\arabic*.,itemsep=2.5mm,leftmargin=*]
  \item Présentez la famille Barker : qui la compose, et quel âge ont les enfants ?
    \lignes{3}
    \corr{Jeremiah Barker, dit Jerry, cocher, propriétaire de son fiacre et de deux chevaux ;
    sa femme Polly, petite, brune, l'air gai ; Harry, douze ans, grand et franc ; Dolly, huit
    ans, le portrait de sa mère. Tous s'aiment beaucoup.}
  \item Que sait-on de Captain, l'autre cheval ?
    \lignes{2}
    \corr{Grand, blanc, un peu fort d'ossature ; vieux maintenant, mais il a dû être
    magnifique. Il porte encore la tête haute. Dans sa jeunesse il a fait la guerre de Crimée,
    chez un officier de cavalerie, et menait le régiment.}
  \item \en{Polly says the horse is too good for a cab, ``if it was not for the broken
    knees''. What does Jerry answer, and what does his answer tell us about him?}
    \lignes{3}
    \corr{Broken knees mean the horse has fallen, and a fall is usually blamed on the horse.
    Jerry says that since nobody can tell whose fault it was, he will give the horse the
    benefit of the doubt. He judges the animal on what he can see, not on its bad reputation
    — the first master in the book to do so.}
  \item \en{Find the only relative clause in the first paragraph. Why is there a comma before
    it? Then rewrite \emph{an animal called ``Captain''} with a relative pronoun.}
    \lignes{3}
    \corr{\emph{Jerry had a cab of his own, and two horses, which he drove and attended to
    himself.} La virgule parce que la relative n'identifie pas les chevaux — il n'en a que
    deux — mais ajoute un renseignement sur eux. \emph{An animal which was called ``Captain''}
    : quand le relatif est sujet et le verbe \emph{be}, on peut supprimer les deux et garder
    le participe seul.}
\end{enumerate}

\note[Un nom qui revient]{\emph{We'll call him `Jack', after the old one} : \emph{to call
somebody after somebody} = donner à quelqu'un le nom d'un autre. Le cheval perd son nom une
fois de plus — c'est le quatrième du livre.}

%% ===================================================================
\section{Version}
\consigne{Traduisez le passage en français. Il précède le texte : le cheval arrive chez son
nouveau maître, le soir même de la foire.}

\begin{version}
My owner pulled up at one of the houses and \gl[to whistle]{whistled}{siffler}. The door flew
open, and a young woman, followed by a little girl and boy, ran out. There was a very lively
greeting as my rider \gl[to dismount]{dismounted}{mettre pied à terre}. ``Now, then, Harry, my
boy, open the gates, and mother will bring us the \gl[a lantern]{lantern}{une lanterne}.'' The
next minute they were all standing round me in a small stable-yard. ``Is he gentle, father?''
``Yes, Dolly, as gentle as your own \gl[a kitten]{kitten}{un chaton}; come and
\gl[to pat]{pat}{caresser, flatter} him.'' At once the little hand was patting about all over
my shoulder without fear. How good it felt!
\end{version}
\source{\emph{Ibid.}, ch.~\textsc{xxxii}.}

\lignes{14}

\corr{\textbf{Proposition de traduction.} Mon propriétaire s'arrêta devant l'une des maisons
et siffla. La porte s'ouvrit à la volée, et une jeune femme, suivie d'une petite fille et d'un
petit garçon, en sortit en courant. Il y eut des retrouvailles très animées tandis que mon
cavalier mettait pied à terre. \og Allons, Harry, mon garçon, ouvre la barrière, et maman nous
apportera la lanterne. \fg{} La minute d'après, ils se tenaient tous autour de moi dans une
petite cour d'écurie. \og Est-il doux, papa ? --- Oui, Dolly, aussi doux que ton propre
chaton ; viens le caresser. \fg{} Aussitôt la petite main me caressa partout l'épaule, sans
crainte. Que c'était bon !}

\note[Choix de traduction 1]{\emph{The door flew open} : \emph{to fly} + adjectif décrit le
mouvement et son résultat d'un seul coup. Le français doit choisir un verbe et ajouter une
manière — \og s'ouvrit à la volée \fg.}

\note[Choix de traduction 2]{\emph{as my rider dismounted} : \emph{as} introduit ici la
simultanéité, non la cause ni la comparaison. \og Tandis que \fg, \og comme \fg.}

\note[Choix de traduction 3]{\emph{How good it felt!} : \emph{to feel} sans complément se
traduit par un impersonnel — \og que c'était bon \fg{} — et non par \og comme il se sentait
bien \fg, qui changerait de sujet.}

%% ===================================================================
\partiel{Lexique à mémoriser}
\begin{multicols}{2}\small
\begin{itemize}[label=--,itemsep=0.8mm,leftmargin=*]
  \item \textbf{a cab} — un fiacre, un taxi
  \item \textbf{a stable} — une écurie
  \item \textbf{a yard} — une cour \textit{(fiche 9)}
  \item \textbf{tidy} — net, rangé
  \item \textbf{fond of} — qui aime beaucoup
  \item \textbf{a lad} — un garçon
  \item \textbf{gentle} — doux \textit{(et non gentil)}
  \item \textbf{a slice} — une tranche
  \item \textbf{a fault} — une faute, la responsabilité
  \item \textbf{a doubt} — un doute ; \textit{le b ne se prononce pas}
  \item \textbf{a knee} — un genou \textit{(fiche 11)}
  \item \textbf{to belong to} — appartenir à
\end{itemize}
\end{multicols}

\end{document}
